% ===== 第 1 章 c01a：Chan 2014，原文 pp.1--5 =====
% [原文 p.1]
\section*{摘要}
本文综述 Strominger--Yau--Zaslow 猜想自 1996 年由 Strominger、Yau 与 Zaslow 在其著名论文 [154] 中提出以来的发展，并讨论它如何引导我们彻底理解镜像对称背后的几何。

\section{SYZ 之前的镜像对称}
1976 年 Yau 证明著名的 Calabi 猜想 [169, 170] 时，大多数数学家------或许包括 Yau 本人在内------都不会想到 Calabi--Yau 流形会在弦论中扮演如此关键而不可替代的角色 [13, 153]；弦论是一门旨在统一广义相对论与量子场论的现代物理理论。人们也不会猜测，几何与物理之间由此产生的互动最终会引出一个重要发现------镜像对称；在二十多年间，这一主题深刻地影响了许多数学分支的发展。

故事始于 20 世纪 80 年代末：Dixon [41] 以及 Lerche、Vafa 与 Warner [110] 提出，理论上两个不同的 Calabi--Yau 流形可以给出完全相同的物理。这一惊人的预言很快被 Greene 与 Plesser [64] 以及 Candelas、Lynker 与 Schimmrigk [14] 所验证：他们各自独立构造出 Hodge 数发生互换的 Calabi--Yau 流形对。我们称这样的流形对为\emph{镜像对}。在 1991 年那篇非凡的论文 [12] 中，Candelas、de la Ossa、Green 与 Parkes 完成了一个更为惊人的计算，由此给出了关于 $\PP^4$ 中五次三维流形上有理曲线数目的预言。数学家对此格外感兴趣，因为这一预言远远超出了当时代数几何学家所能达到的范围。

这催生了 Gromov--Witten 理论等若干重要课题，并最终由 Givental [60, 61] 与 Lian--Liu--Yau [116, 117, 118, 119] 各自独立证明镜像定理而达到高潮；这些证明特别验证了 Candelas 等人关于五次三维流形的预言。这无疑是一项辉煌的成就。然而，所有这些证明都依赖于包含该 Calabi--Yau 流形的外围环面簇的几何；尤其是，它们并未提供一种内蕴的方式来理解镜像对称的几何。

% [原文 p.2]
镜像对称的第一个内蕴数学表述，是 Kontsevich 在其 1994 年 ICM 报告 [103] 中提出的同调镜像对称（HMS）猜想。在弦论中，一个 Calabi--Yau 流形 $X$ 确定两个拓扑弦理论：A-模型与 B-模型，它们分别依赖于 $X$ 的辛几何与复几何 [161, 167]。从这个观点看，镜像对称预言：若 $X$ 与 $\check X$ 是一对镜像 Calabi--Yau 流形，则 $X$ 的 A-模型与 $\check X$ 的 B-模型同构，反之亦然。上文关于五次三维流形的计数预言，只是这一更大图景的众多有趣表现之一。

Kontsevich 的 HMS 猜想把镜像对称简洁地表述为：$X$ 中拉格朗日子流形的 Fukaya 范畴（A-模型）与镜像 $\check X$ 上凝聚层的导出范畴（B-模型）之间的等价。他的猜想既深刻又优美，并且被指望能推出镜像对称的计数预言。然而，它既未指出如何找到这样的等价，也未告诉我们如何构造给定 Calabi--Yau 流形的镜像。

\section{SYZ 猜想的表述}
1996 年夏，Strominger、Yau 与 Zaslow [154] 提出了一个开创性的设想，首次为镜像对称给出几何解释：

\begin{conjecture}{1.2.1}[SYZ 猜想，见 {[}154{]}]
设 $X$ 与 $\check X$ 是互为镜像的 Calabi--Yau 流形。则
\begin{enumerate}
  \item $X$ 与 $\check X$ 都容许带截面的特殊拉格朗日环面纤维化 $\mu : X \to B$ 与 $\check\mu : \check X \to B$，且底空间相同：
  \begin{center}
  \begin{tikzcd}
    X \arrow[d, "\mu"] & \check X \arrow[d, "\check\mu"] \\
    B \arrow[r, equal] & B
  \end{tikzcd}
  \end{center}
  \item 纤维化 $\mu : X \to B$ 与 $\check\mu : \check X \to B$ 在逐纤维对偶的意义下互为对偶：即若 $b \in B$ 上方的纤维 $\mu^{-1}(b) \subset X$ 与 $\check\mu^{-1}(b) \subset \check X$ 非奇异，则它们是互为对偶的环面；并且
  \item 存在逐纤维的 Fourier（--Mukai）型变换，它们负责在 $X$ 上的辛几何数据（相应地，复几何数据）与 $\check X$ 上的复几何数据（相应地，辛几何数据）之间进行互换。
\end{enumerate}
\end{conjecture}

一言以蔽之，这一猜想是说：神秘的镜像现象不过是 Fourier 变换！这个非凡而影响深远的猜想不仅为镜像对称提供了一个优美的几何解释，还暗示给定 Calabi--Yau 流形 $X$ 的镜像伙伴可以通过对 $X$ 上的特殊拉格朗日环面纤维化逐纤维取对偶来构造（也就是 T 对偶）。它立刻吸引了数学家和物理学家的广泛关注，并促成了大量研究：或者求解这一猜想，或者应用它来理解镜像对称背后的几何。

在继续之前，让我们简要回顾 SYZ 猜想背后的启发式论证。首先，弦论的一个关键特征是 Dirichlet 膜（即 D-膜）的存在。物理上的论证表明：B-模型中的 D-膜（或简称 B-膜）是复子簇上的凝聚层，而 A-模型中的 D-膜（或简称 A-膜）是配备平坦联络的特殊拉格朗日子流形。由于镜像对称预言 $X$ 的 A-模型与 $\check X$ 的 B-模型同构，$X$ 上 A-膜的模空间应当与 $\check X$ 上镜像 B-膜的模空间等同。

% [原文 p.3]
现在，$\check X$ 上的点当然可以看作 B-膜。又因为 $\check X$ 本身就是这些 B-膜的模空间，所以它应当与 $X$ 上某些 A-膜 $(L,\nabla)$ 的模空间等同，其中 $L \subset X$ 是特殊拉格朗日子流形，$\nabla$ 是 $L$ 上的平坦 $U(1)$-联络。此外，既然 $\check X$ 被它的点所扫成，$X$ 也应当被这些特殊拉格朗日子流形 $L$ 所扫成。由 McLean 定理 [127]，特殊拉格朗日子流形 $L \subset X$ 的模空间是无障碍的，并以 $H^1(L;\R)$ 为模型；而 $L$ 上平坦 $U(1)$-联络（模去规范等价）的模空间由 $H^1(L;\R)/H^1(L;\Z)$ 给出。因此，为使维数匹配，我们应有 $\dim H^1(L;\R) = \dim_\C \check X = n$。于是 $X$ 应容许一个特殊拉格朗日环面纤维化
\[
  \mu : X \to B .
\]
此外，流形 $\check X$ 本身可以看成一个 B-膜，其模空间是单点集，并且它与 $\check X$ 中每个点恰好相交一次；因此对应的 A-膜应给出 $\mu$ 的一个特殊拉格朗日截面 $\sigma$，满足 $H^1(\sigma;\R)=0$。特别地，底空间 $B$ 的第一 Betti 数应为 $b_1 = 0$。

把同样的论证应用于 $\check X$，得到带截面的特殊拉格朗日环面纤维化
\[
  \check\mu : \check X \to \check B .
\]
现在，对于环面纤维 $L_b := \mu^{-1}(b) \subset X$，其对偶 $L_b^\vee$ 可以看作 $L$ 上平坦 $U(1)$-联络的模空间；在镜像对称下，这些联络对应于 $\check X$ 中的点。这表明 $L_b^\vee$ 是 $\check X$ 的子流形。经过更精细的论证可以看出，$L_b^\vee$ 实际上可以与 $\check\mu$ 的一个特殊拉格朗日环面纤维等同，从而推出 $\mu$ 与 $\check\mu$ 是同一底空间上的纤维化，并且逐纤维互为对偶。

注意，我们有一个变换，把 $X$ 中的特殊拉格朗日环面纤维（A-膜）映为 $\check X$ 中的点（B-膜）。这是逐纤维的 Fourier（--Mukai）型变换的一个实例。更一般地，应当存在几何的 Fourier 变换，把 $X$ 上的辛几何数据映为 $\check X$ 上的复几何数据。我们称这些变换为 \emph{SYZ 变换}。在原始的 SYZ 论文 [154] 中，人们推断：镜像 $\check X$ 上 Ricci 平坦度量的行为，应当由瞬子修正的贡献而不同于此前 Greene、Shapere、Vafa 与 Yau 在一篇重要论文 [65] 中构造的半平坦 Calabi--Yau 度量。正如我们将看到的，研究镜像对称的一个关键步骤就是理解这些修正，它们应当来自 SYZ 变换的高阶 Fourier 模。

\section{半平坦 SYZ}
当特殊拉格朗日环面纤维化不含任何奇异纤维时，SYZ 图景特别优美。McLean 的经典结果 [127] 给出特殊拉格朗日环面纤维化的底流形 $B$ 上两个自然定义的\emph{整仿射结构}\footnote{流形上的整仿射结构是一个图册，其转移映射全为整仿射线性变换。}：辛仿射结构与复仿射结构；镜像对称可以通过这些结构得到简洁的解释。更具体地说，设 $v$ 是纤维 $L_b := \mu^{-1}(b)$ 的一个法向量场，它确定一个 1-形式 $\alpha := -\iota_v \omega \in \Omega^1(L_b;\R)$ 与一个 $(n-1)$-形式 $\beta := \iota_v \operatorname{Im}\Omega \in \Omega^{n-1}(L_b;\R)$，其中 $\omega$ 与 $\Omega$ 分别是 $X$ 上的 K\"ahler 形式与全纯体积形式。McLean [127] 证明了对应的形变是特殊拉格朗日的，当且仅当 $\alpha$ 与 $\beta$ 都闭。利用 $\alpha$ 的上同调类把 $TB$ 与 $H^1(L_b;\R)$ 等同起来，我们

% [原文 p.4]
就得到 $B$ 上的辛仿射结构；而利用 $\beta$ 的上同调类把 $TB$ 与 $H^{n-1}(L_b;\R)$ 等同起来，就得到 $B$ 上的复仿射结构。我们还有由下式定义的 \emph{McLean 度量}
\[
  g(v_1,v_2) := -\int_{L_b} \iota_{v_1}\omega \wedge \iota_{v_2}\operatorname{Im}\Omega .
\]
在他那篇富有启发性的论文 [87] 中，Hitchin 解释了这些结构如何通过 \emph{Legendre 变换}相互联系。若用 $x_1,\dots,x_n$ 记 $B$ 上关于辛仿射结构的局部仿射坐标，则局部地，McLean 度量可以写成 $B$ 上一个凸函数 $\phi$ 的 Hessian，即
\[
  g\!\left(\frac{\partial}{\partial x_i},\frac{\partial}{\partial x_j}\right)
  = \frac{\partial^2\phi}{\partial x_i\partial x_j} .
\]
此外，令 $\check x_i := \partial\phi/\partial x_i$（$i=1,\dots,n$），恰好给出 $B$ 上关于复仿射结构的局部仿射坐标；并且若
\[
  \check\phi := \sum_{i=1}^{n} \check x_i x_i - \phi(x_1,\dots,x_n)
\]
是 $\phi$ 的 Legendre 变换，则我们有 $x_i = \partial\check\phi/\partial\check x_i$ 以及
\[
  g\!\left(\frac{\partial}{\partial\check x_i},\frac{\partial}{\partial\check x_j}\right)
  = \frac{\partial^2\check\phi}{\partial\check x_i\partial\check x_j} .
\]
若进一步假设纤维化 $\mu : X \to B$ 容许一个拉格朗日截面，则由 Duistermaat 的一个定理 [42] 可知存在整体作用量--角变量坐标，使得我们可以写出
\[
  X = T^*B/\Lambda^\vee ,
\]
其中格 $\Lambda^\vee \subset T^*B$ 局部由 $dx_1,\dots,dx_n$ 生成，并且 $\omega$ 可以等同于 $T^*B/\Lambda^\vee$ 上的典范辛形式
\[
  \omega = \sum_{i=1}^{n} dx_i \wedge du_i .
\]
这里 $u_1,\dots,u_n$ 是 $T^*B$ 上的纤维坐标。

在这种情况下，$X$ 的镜像简单地由
\[
  \check X := TB/\Lambda
\]
给出，其中格 $\Lambda \subset TB$ 局部由 $\partial/\partial x_1,\dots,\partial/\partial x_n$ 生成。商流形 $\check X$ 具有自然的复结构，其全纯坐标由 $z_i := \exp(x_i + \ii y_i)$ 给出，其中 $y_1,\dots,y_n$ 是 $TB$ 上与 $u_1,\dots,u_n$ 对偶的纤维坐标。这样就把 $X$ 的镜像构造为一个带有处处非零全纯体积形式的复流形
\[
  \check\Omega := d\log z_1 \wedge \cdots \wedge d\log z_n .
\]
此外，存在一个显式的逐纤维 Fourier 型变换，我们称之为\emph{半平坦 SYZ 变换} $\mathcal{F}^{\text{semi-flat}}$，它把 $e^{\ii\omega}$ 变为 $\check\Omega$；详见 [23, 第 2 节]。

现在，若我们改用 $B$ 上的复仿射结构，就得到 $\check X$ 上一个与其复结构相容的辛结构，使得镜像 $\check X$ 成为一个 K\"ahler 流形。进一步，若上述函数 $\phi$ 满足实 Monge--Amp\`ere 方程
\[
  \det\left(\frac{\partial^2\phi}{\partial x_i\partial x_j}\right) = \text{常数},
\]
则我们在 $X$ 与其镜像 $\check X$ 上都得到 $T^n$-不变的 Ricci 平坦度量。$B$ 上诱导的度量称为 \emph{Monge--Amp\`ere 度量}，$B$ 称为 \emph{Monge--Amp\`ere 流形}。

% [原文 p.5]
这把镜像对称与实 Monge--Amp\`ere 方程以及仿射 K\"ahler 几何的研究联系起来；在镜像对称被发现之前，Cheng 与 Yau 已经在这一方向上做出了重要贡献 [32, 33, 34]。在带奇点的仿射流形上构造 Monge--Amp\`ere 度量，此后一直是仿射几何与 SYZ 猜想研究中的一个重要问题。Loftin、Yau 与 Zaslow 的高度非平凡的工作 [122, 123] 在“Y”型顶点附近构造了这样的度量；“Y”型顶点是三维情形中一种典型的奇点类型。但除此之外，人们所知甚少。

所以，在半平坦情形，SYZ 猜想确实为镜像对称描绘了一幅吸引人的图景；关于半平坦 SYZ 镜像对称的更多细节由 Leung 在 [111] 中给出。不幸的是，这幅优美的图景只能在所有瞬子修正都被抑制的大复结构/体积极限下成立。离开这些极限，特殊拉格朗日纤维化将出现奇异纤维，镜像不再能简单地通过对一个纤维化取对偶而得到。

\section{构造 SYZ 纤维化}
1996 年 SYZ 猜想提出之后，人们立刻投入了大量精力，试图在 Calabi--Yau 流形上构造特殊拉格朗日环面纤维化（即 SYZ 纤维化）\footnote{在 [67, 68] 中，Gross 并未构造这样的纤维化，而是假设其存在，并推出镜像对称所预言的若干有趣推论。}。Zharkov [174] 第一个在光滑射影环面簇 $P_\Delta$ 中的 Calabi--Yau 超曲面构造了拓扑环面纤维化。这包括了五次三维流形这一重要例子。Zharkov 的构造方式是：把 $P_\Delta$ 上矩映射在矩多面体边界 $\partial\Delta$ 上的限制，形变到附近一个光滑的 Calabi--Yau 超曲面上。

利用类似的想法并引入梯度--Hamilton 流，W.-D. Ruan 在一系列论文 [142, 143, 144] 中在五次三维流形上构造了拉格朗日环面纤维化。他还完成了这些纤维化单值群的一项重要计算；此后 Gross [70] 利用这一计算给出了包括五次三维流形在内的 Calabi--Yau 流形的拓扑版 SYZ 镜像对称。Mikhalkin [128] 也有一项相关工作：他运用热带几何的工具，在环面簇的超曲面上生成光滑环面纤维化。

一般而言，构造特殊拉格朗日纤维化、甚至构造特殊拉格朗日子流形，都是非常困难的问题。构造特殊拉格朗日子的一个有希望的方法是使用平均曲率流。Thomas [156] 为 Calabi--Yau 流形中 Maslov 指标为零的拉格朗日子流形类建立了一个稳定性概念；它应当镜像于全纯向量丛的稳定性。在 [157] 中，Thomas 与 Yau 猜想：一个 Hamiltonian 同痕类中存在唯一的特殊拉格朗日代表元，当且仅当该类是稳定的；并且这样的代表元可以通过平均曲率流得到，而平均曲率流应当长时间存在。他们还为特殊拉格朗日子流形提出了 Jordan--H\"older 型分解，并将其与平均曲率流中奇点的形成联系起来。他们的方案与猜想对 Calabi--Yau 几何与 SYZ 猜想的发展产生了重大影响。这一领域已取得许多进展 [29, 30, 96, 109, 131, 132, 134, 148, 149, 150, 162, 163, 164, 165]；更详细的内容可参见出色的综述文章 [166] 与 [133] 及其中的参考文献。
